% !TeX root = ../colloquium-2.tex

% Формулировки следствий теоремы Лагранжа.
% Формулировка теоремы Коши.

\phantomsection\label{sec_Definition13}
\begin{definition}[Формулировки следствий теоремы Лагранжа.
                   Формулировка теоремы Коши.]
    \subsection*{Формулировки следствий теоремы Лагранжа.}
    \begin{note}[amber]
        \subsubsection*{Следствие 1.}
        \begin{note}[purple]
            Пусть определена и дифференцируема на интервале $(a, \; b)$
            и $f'(x) = 0 \; \forall x \in (a, \; b)$. Тогда функция $f$ постоянна на $(a, \; b)$. 
        \end{note}

        \subsubsection*{Следствие 2.}
        \begin{note}[purple]
            Пусть функция $f$ определена и дифференцируема на интервале $(a, \; b)$
            \begin{enumerate}
                \item Функция $f$ не убывает на этом интервале тогда и только тогда, 
                когда $$ f'(x) \geq 0 \; \forall x \in (a, \; b) $$
                \item Функция $f$ не возрастает на этом интервале тогда и только тогда, 
                когда $$ f'(x) \leq 0 \; \forall x \in (a, \; b) $$
            \end{enumerate}
        \end{note}
    \end{note}

    \subsection*{Формулировка теоремы Коши.}
    \begin{note}[purple]
        Пусть функции $f$ и $g$:
        \begin{enumerate}
            \item определены и непрерывны на отрезке $[a,\; b]$
            \item дифференцируемы на интервале $(a, \; b)$
            \item $g'(x) \neq 0 \; \forall x \in (a, \; b)$
        \end{enumerate}
        Тогда существует такая точка $c \in (a, \; b)$, что 
        $$ \frac{f(b) - f(a)}{g(b) - g(a)} = \frac{f'(c)}{g'(c)} $$
    \end{note}
\end{definition}